\documentclass[12pt, letterpaper]{article}

%-------Packages---------
\usepackage{amssymb,amsfonts}
\usepackage{mathptmx}
\usepackage{amsthm}
\usepackage{graphicx}
\usepackage[all,arc]{xy}
\usepackage{enumerate}
\usepackage{bbm}
\usepackage{dsfont}
\usepackage{mathrsfs}
\usepackage{tikz-cd}
\usepackage{setspace}
\usepackage{import}
\usepackage[utf8]{inputenc}
\usepackage{hyperref}
\usepackage{tikz}
\usepackage{float}
\usepackage[dvipsnames]{xcolor}
\usepackage[nocheck]{fancyhdr}
\usepackage[nointegrals]{wasysym}

\usepackage[left=1in,top=1in,right=1in,bottom=1in]{geometry}

\usepackage{mathtools}
\usepackage{setspace}
\usepackage{cleveref}
\usepackage{multicol}
\usepackage{mathpazo}
\usepackage{tcolorbox}
\tcbuselibrary{theorems,skins,breakable}
\usepackage{xparse}

\usepackage{xcolor, ifthen, tcolorbox}
\tcbuselibrary{theorems,skins,breakable}
% Theorem System
% The following environments are provided:
%   New Problem:        \begin{problem} ...
%   New Question Header:   \qh (then followed by subproblems)
%   Subproblem:     \begin{subproblem} ...
%   Problem Proof:  \begin{problemproof} ...
%   Proof:          \\begin{pf}{color} ...
%   Remark:         \rmk (sentence), \rmkb (block)

% Define custom colors
\definecolor{thmscol}{HTML}{F4565B} %For Problems
\definecolor{lemscol}{HTML}{FE844C} %For Lemmes
\definecolor{prpscol}{HTML}{00A7EF} %For proposition
\definecolor{corscol}{HTML}{23DAFF} %For corrolaries
\definecolor{rmkscol}{HTML}{6DEEC5} %For remarks
\definecolor{defscol}{HTML}{18C991} %For definitions
\definecolor{exmscol}{HTML}{7DB800} %For examples
\definecolor{clmscol}{HTML}{165c58} %For claims

%Change between dark(black)/light(white) mode
\newcommand{\colormode}{white}
\ifthenelse{\equal{\colormode}{white}}
      {\newcommand{\TextColor}{black}}
      {\newcommand{\TextColor}{white}}
\newcommand{\pbcolormain}{thmscol!80!\colormode}
\newcommand{\pbcolorsecond}{thmscol!38!\colormode}


% ============================
% Problem Counter
% ============================
\newcounter{subproblemcounter}
\newcounter{problemcounter}

\newcommand{\q}{
    \setcounter{subproblemcounter}{0}%
    \stepcounter{problemcounter} % Increment problem counter
    \color{\TextColor}Problem \arabic{problemcounter}: % Display the problem counter
}

\newcommand{\stepsub}{
    \stepcounter{subproblemcounter}%
    \color{\TextColor}(\alph{subproblemcounter})
}
% ============================
% Problem
% ============================
\newtcolorbox{pb}[1]{
  enhanced,
  frame hidden,
  titlerule=0mm,
  toptitle=1mm,
  bottomtitle=1mm,
  fonttitle=\bfseries\large,
  coltitle=black,
  colbacktitle=\pbcolormain,
  colback=\pbcolorsecond,
  title={\q #1},
  coltext=\TextColor
}

\newtcolorbox{spb}[1]{
  enhanced,
  frame hidden,
  titlerule=0mm,
  toptitle=1mm,
  bottomtitle=1mm,
  fonttitle=\bfseries\large,
  coltitle=black,
  colbacktitle=\pbcolormain,
  colback=\pbcolorsecond,
  title={\stepsub #1},
  coltext=\TextColor,
  attach boxed title to top left={xshift=3mm,yshift=-2mm},
  boxed title style={
    colback=\pbcolormain,
    colframe=\pbcolormain,
  }
}

\newtcolorbox{pbh}[1]{
    enhanced,
    frame hidden,
    titlerule=0mm,
    toptitle=1mm,
    bottomtitle=1mm,
    fonttitle=\bfseries\large,
    coltitle=black,
    colbacktitle=\pbcolormain,
    colback=\pbcolorsecond,
    title={\q #1},
    boxrule=0pt,
    interior hidden,
    attach boxed title to top left={xshift=3mm,yshift=-2mm},
    boxed title style={
        colback=\pbcolormain,
        colframe=\pbcolormain,
    }
}

\newenvironment{problem}[1]{%
  \newpage
  \begin{pb}{#1}
    }{
  \end{pb}
}

\newenvironment{subproblem}[1]{%
  \begin{spb}{#1}
    }{
  \end{spb}
}

\newenvironment{problemproof}{%
    \begin{pf}{\pbcolormain}
        }{
    \end{pf}
}

\newcommand{\qh}[1][]{
    \newpage
    \begin{pbh}{#1}\end{pbh} % Display the problem counter
}

% ============================
% Claim
% ============================

\newtcbtheorem[number within=problemcounter]{claim}{Claim}
{
    enhanced,
    frame hidden,
    titlerule=0mm,
    toptitle=1mm,
    bottomtitle=1mm,
    fonttitle=\bfseries\large,
    coltitle=black,
    colbacktitle=clmscol!50!\colormode,
    colback=clmscol!20!\colormode,
}{clm}

\NewDocumentCommand{\clm}{m+m}{
    \begin{claim*}{#1}{}
        #2
    \end{claim*}
}

\newenvironment{clmpf}{
	{\noindent{\it \textbf{Proof.}}
	\setlength{\parindent}{0pt}
	}
	\tcolorbox[blanker,breakable,left=5mm,parbox=false,
    before upper={\parindent15pt},
    after skip=10pt,
	borderline west={1mm}{0pt}{clmscol!50!\colormode}]
}{
    \textcolor{clmscol!50!\colormode}{\hbox{}\nobreak\hfill$\blacksquare$}
    \endtcolorbox
}

\NewDocumentCommand{\clmp}{m+m+m}{
    \begin{claim*}{#1}{}
        #2
    \end{claim*}

    \begin{clmpf}
        #3
    \end{clmpf}
}


% ============================
% Proof
% ============================

\newenvironment{pf}[1]{%
  \def\pfcolor{#1}% Store the color in a macro
  \noindent{\it \textbf{Proof.}}%
  \tcolorbox[blanker,breakable,left=5mm,parbox=false,%
    before upper={\parindent15pt},%
    after skip=10pt,%
    borderline west={1mm}{0pt}{\pfcolor},%
    coltext=\TextColor
  ]%
  \setlength{\parindent}{0pt}%
}{%
  \textcolor{\pfcolor}{\hbox{}\nobreak\hfill$\blacksquare$}%
  \endtcolorbox%
}

\newtcolorbox{solution}{
  blanker,
  breakable,
  left=5mm,
  parbox=false,%
  before upper={\parindent15pt},%
  after skip=10pt,%
  borderline west={1mm}{0pt}{\pbcolormain},%
  coltext=\TextColor
}


% ============================
% Remark
% ============================


\NewDocumentCommand{\rmk}{+m}{
    {\it \color{rmkscol!100!white}#1}
}

\newenvironment{remark}{
    \par
    \vspace{5pt}
    \begin{minipage}{\textwidth}
        {\par\noindent{\textbf{Remark.}}}
        \tcolorbox[blanker,breakable,left=5mm,
        before skip=10pt,after skip=10pt,
        borderline west={1mm}{0pt}{rmkscol!80!\colormode},
        coltext=\TextColor
        ]
}{
        \endtcolorbox
    \end{minipage}
    \vspace{5pt}
}

\NewDocumentCommand{\rmkb}{+m}{
    \begin{remark}
        #1
    \end{remark}
}


% Define a new command for problems
\renewcommand{\headrule}{\color{\TextColor}\hrulefill}
\newcommand{\C}{\mathbb{C}}
\newcommand{\R}{\mathbb{R}}
\newcommand{\Q}{\mathbb{Q}}
\newcommand{\Z}{\mathbb{Z}}
\newcommand{\N}{\mathbb{N}}
\renewcommand{\P}{\mathbb{P}}
\newcommand{\F}{\mathbb{F}}
\newcommand{\T}{\mathcal{T}}
\newcommand{\contr}{\Rightarrow \Leftarrow}
\newcommand{\s}{\(\,\)}
\renewcommand{\t}[1]{\text{#1}}
\newcommand{\ang}[1]{\left\langle #1 \right\rangle}
\newcommand{\coker}{\mathrm{coker}}

% Meta data
\setstretch{1.2}

\begin{document}
\color{\TextColor}
\pagecolor{\colormode}
\setlength{\parindent}{0pt}
\pagestyle{fancy}
\fancyhf{}
\fancyhead[LOE]{\color{\TextColor}\slshape{\textbf{Vincent Ferrara}}}
\fancyhead[ROE]{\color{\TextColor}\thepage}

\begin{problem}{}
    A space $X$ is \emph{locally metrizable} if each point in $X$ has a
    neighborhood which is metrizable in the subspace topology. Prove
    that if $X$ is a locally metrizable compact Hausdorff space, then it
    is metrizable. \textit{Hint:} you should use the Urysohn metrization
    theorem; for a further hint, see Exercise 7 in \S 34 of Munkres.
\end{problem}
\begin{problemproof}
    Assume $X$ is locally metrizable compact Hausdorff space. Thus, by homework $X$ is normal since it is compact and Hausdorff. Therefore, $X$ is also regular.

    Let $x \in X$. Let $U_x$ be a neighborhood of $x$ s.t. it is metrizable. Then define $B_x=\{U_x \cap B(x,1/n) \mid n \in \N\} \cup \{U_x \cap B(x,n) \mid n \in \N\}$.

    Consider $\{U_x\}_{x \in X}$. This is an open cover of $X$. Thus, there is a finite subcover, $U_{x_1},\dots,U_{x_n}$ since $X$ is compact.

    Since $U_x$ is open in $X$ and $B(x,\epsilon)$ is open in $U_x$. This implies that $B(x,\epsilon)=U_x \cap O$ for $O$ open in $X$. Thus, $U_x \cap B(x,\epsilon)=U_x \cap O$ is open in $X$.

    Define $\beta$ as all finite intersections of $B=B_{x_1} \cup \dots \cup B_{x_n}$

    Basis:

    Let $x \in X$. Then $x \in U_{x_i}$ for some $i$. Thus, let $d$ be the metric for $U_{x_i}$. Thus, $d(x,x_i)=a$ for some $a \in \R$. Thus, fix $N \in \N$ s.t. $N > a$. Thus, $x \in B(x_i,N) \cap U_{x_i} \in \beta$.

    Let $x \in B_1 \cap B_2$ s.t. $B_1 \cap B_2 \in \beta$. Thus, $B_1 \cap B_2 \in \beta$. Since $B_1$ is a finite intersection of things in $B$ and same with $B_2$, so their intersection is also a finite intersection.

    Thus, $x \in B_1 \cap B_2 \subset B_1 \cap B_2$.

    Thus, $\beta$ is a countable basis for $X$.

    Thus, by the Urysohn metrization theorem we know that $X$ is metrizable.
\end{problemproof}

\begin{problem}{}
    Let $X$ be a topological space and let $Y$ be a metric space. A
  function $f:X \to Y$ is \emph{bounded} if there exists some $R > 0$
  such that, for any $x_1, x_2 \in X$, we have
  $d(f(x_1), f(x_2)) \le R$. The set of all bounded functions
  $f:X \to Y$ is denoted $\mathcal{B}(X,Y)$.
\end{problem}
\begin{subproblem}{}
    Is $\mathcal{B}(X, Y)$ necessarily closed in the topology of
    pointwise convergence on $Y^X$?
\end{subproblem}
\begin{problemproof}
    No consider $\R^\R$
    \begin{equation*}
        f_n(x)=\begin{cases*}
            x \quad -n \leq x \leq n,\\
            -n \quad -n < x,\\
            n \quad n > x.
        \end{cases*}
    \end{equation*}
    $f_n \in \mathcal{B}(\R,\R)$ but $f_n \to f(x)=x$ since fix $\epsilon > 0$, and $a \in \R$.

    Fix $N \in \N$ big enough s.t. $|a| < N$. Thus, $d(a,f(a))=0 < \epsilon$. Thus, $f_x(x) \to f(x)$ for all $x \in \R$. Thus, $f_x \to f$, but $f$ is not bounded, so $\mathcal{B}(X,Y)$ is not closed since if it was closed it would contain all of its limit points.
\end{problemproof}

\begin{subproblem}{}
    Is $\mathcal{B}(X, Y)$ necessarily closed in the topology of
    uniform convergence on $Y^X$?
\end{subproblem}
\begin{problemproof}
    Let $(f_n)$ be a convergent sequence of functions in $\mathcal{B}(X,Y)$ s.t. $(f_n) \to f$. Fix $\epsilon > 0$. Thus, there is an $N \in \N$ s.t. $d(f_n(x),f(x)) < \epsilon$ for all $n \geq N$ and for all $x \in X$.

    Thus, consider $d(f(x),f(y))$ s.t. $x,y \in X$. By the triangle inequality

    \begin{align*}
        d(f(x),f(y)) &\leq d(f(x),f_N(x)) + d(f(y),f_N(x)) \\
        &\leq d(f(x),f_n(x)) + d(f(y),f_N(y)) + d(f_N(x),f_N(y)) \\ &\leq 2\epsilon + R_N
    \end{align*}
    for some $R_N \in \R$ since $f_N$ is bounded.

    Thus, $f$ is bounded. Since $\mathcal{B}(X,Y)$ contains all of its limit points this implies that it is closed.
\end{problemproof}

\newcommand{\mc}[1]{\mathcal{#1}}
\begin{problem}{}
    Let $X$ and $Y$ be metric spaces; recall that $\mc{C}(X, Y)$ denotes
  the set of continuous functions $X \to Y$. There is a map
  $e:X \times \mc{C}(X, Y) \to Y$ defined by the rule:
  \[
    e(x, f) = f(x).
  \]
  This map is called the \emph{evaluation map.} Prove that, if
  $\mc{C}(X, Y)$ has the topology of uniform convergence, then $e$ is
  a continuous map.
\end{problem}
\begin{problemproof}
    Let $U \subset Y$ be open. Let $(x_0,f_0) \in e^{-1}(U)$.

    Thus, $f_0(x_0) \in U$. Therefore, there exists $\epsilon > 0$ s.t. $B(f_0(x_0), \epsilon) \subset U$.

    Since $f_0$ is continuous, there exists an open neighborhood $V \subset X$ of $x_0$ s.t. 

    $f(V) \in B(f_0(x),\epsilon/2)$.

    Consider the open ball $B(f_0,\epsilon/2) \subset \mathcal{C}(X,Y)$ this implies that $d_Y(f(x),f_0(x)) \leq \epsilon / 2$ for all $x \in X$.

    Thus, for $(x,f) \in V \times W$

    \[d_Y(f(x),f_0(x)) \leq d_Y(f(x),f_0(x)) + d_Y(f_0(x),f_0(x_0)) \leq \epsilon\]

    Thus, $f(x) \in B(f_0(x_0)) \subset U$.

    Thus, $(x_0,f_0) \in V \times W \subset e^{-1}(U)$. Therefore, $e^{-1}(U)$ is open.
\end{problemproof}

\begin{problem}{}
    Let $X$ be a space, and let $F$ be the set of all functions
  $X \to \R$. Suppose that the uniform topology and the point-open
  topology on $F$ agree. What must be true about $X$?
\end{problem}
\begin{problemproof}
    Assume for contradiction that $X$ is infinite. Thus, let $\{x_n\}$ be a countably infinite subset of $X$.

    Define $f_n(x)=1$ if $x=x_n$ or $0$ otherwise.

    Fix $\epsilon >0$ and $x \in X$.

    If $x \notin \{x_n\}$. Then $d(f(x),0)=0 < \epsilon$ for all $n \in \N$.

    If $x \in \{x_n\}$ then $x=x_i$. Fix $N \in \N$ s.t. $N > i$. Thus, $d(f_n(x),0)=0 < \epsilon$ for all $n \geq N$.

    Thus, $(f_n) \to 0$.

    However, $\sup_{x \in X}d(f_n(x)-0)=1$ for all $n \in \N$. Thus, $f_n$ does not converge uniformly to $0$.

    Thus, this is a contradiction thus $X$ is finite.
\end{problemproof}

\begin{problem}{}
    Consider the sequence of functions $f_n:(-1, 1) \to \R$ defined by
  \[
    f_n(x) = \sum_{k=1}^n kx^k.
  \]
\end{problem}

\begin{subproblem}{}
    Prove that $f_n$ converges pointwise, but \emph{not}
    uniformly.
\end{subproblem}
\begin{problemproof}
    Consider the power seires

    \[\sum_{k=1}^\infty x^k = \dfrac{1}{1-x} \quad \t{for } |x| < 1\]

    Thus, taking the derivative and multiplying by $x$ we get.

    \[\sum_{k=1}^\infty kx^k = \dfrac{x}{(1-x)^2} \quad \t{for } |x| < 1\]

    Thus, let $f(x)=\dfrac{x}{(1-x)^2}$.

    Fix $\epsilon > 0$ and $x \in (-1,1)$

    Thus, \[|f(x)-f_n(x)|=\lvert \sum_{k=n+1}^\infty kx^k \rvert \leq \sum_{k=n+1}^\infty k|x|^k = \sum_{j=0}^\infty (n+1+j)|x|^{n+1+j}\]
    \[=|x|^{n+1}((n+1) \sum_{j=0}^\infty|x|^j + \sum_{j=0}^\infty j|x|^j)=|x|^{n+1}\left(\dfrac{n+1}{1-|x|} + \dfrac{|x|}{(1-|x|)^2}\right)\]

    Since $(|x|^{n+1}(n+1)) \to 0$ and $(C|x|^{n+1}) \to 0$ for $C=|x|/(1-|x|)^2$. This is because $|x| < 1$.

    Thus, there is an $N_1,N_2 \in \N$ s.t. $|x|^{n+1}(n+1) < \epsilon/2$ and $C|x|^{n+1} < \epsilon/2$.

    For the last part use $N=\max\{N_1.N_2\}$

    Thus,

    \[|x|^{n+1}\left(\dfrac{n+1}{1-|x|} + \dfrac{|x|}{(1-|x|)^2}\right) < \epsilon\]

    Therefore, $f_n(x)$ converges pointwise to $f$ on $(-1,1)$

    Suppose $f_n$ converges uniformly to $f$ on $(-1,1)$. The for $\epsilon =1$ there exists an $N \in \N$ s.t. for all $n \geq N$ and all $x \in (-1,1)$
    \[|f_n(x)-f(x)| <1\]

    Shown above we know that $(n+1)x^{n+1} \leq |f_n(x)-f(x)|=R(x) < 1$.

    Consider $x_n=1-(1/n)$

    Thus, $(n+1)(1-1/n)^{n+1} \leq R(x_n)$.

    Thus, as $n \to \infty$ shown in analysis $(1-1/n)^{n+1} \to e^{-1}$.

    Therefore, there exists an $N_1 \in \N$ s.t. for all $n \geq N_1$ 
    
    $-1/2e < (1-1/n)^{n+1}-1/e < 1/2e \implies 1/2e < (1-1/n)^{n+1}$.

    Thus, for such $n \geq N_1$, $R(x_n) \geq (n+1)/(2e)$.

    Also for $(n+1)/(2e)$ consider $N_2=3$ s.t. for all $n \geq N_2$ $(n+1)/2e > 1$.

    Thus, for $\max{N,N_1,N_2}$ we have $R(x_n) \geq (n+1)/2e > 1$. This is a contradiction thus, $f_n$ does not uniformly converge.
\end{problemproof}

\begin{subproblem}{}
    Prove that $f_n$ converges in the
    topology of compact convergence, so its limit is a continuous
    function.
\end{subproblem}
\begin{problemproof}
    Let $K \subset (-1,1)$ since $(-1,1)$ is a subspace of $\R$. This implies that $K$ is closed and bounded. Thus, there exists an $r \in (0,1)$ s.t. $K \subset [-r,r]$.

    Let $M_k=kr^k$.

    Define a new sequence $g_n(x)=nx^n$.

    Thus, for $x \in K$ we have $|g_n(x)| \leq n|x|^n \leq M_n$.

    Since $\sum_{k=1}^\infty M_k$ converges by standard power series rules which is shown above since $|r|<1$. We get that by the Weierstrass M-test that $\sum_{k=1}^\infty g_n(x)$ converges uniformly on $[-r,r]$.

    Thus, this implies that $f_n(x)$ converges uniformly to on $K \subset [-r,r]$.

    Also, since $(-1,1)$ is a subspace of $\R$ it is metrizable thus first countable. Thus, it is compactly generated. Also since $\R$ is metrizable this implies that $\mathcal{C}((-1,1),\R)$ is closed in the topology of compact convergence which implies that the limit of $f_n$ is continuous since each $f_n$ is continuous.
\end{problemproof}
\end{document}

